\calculussheet{sheet:p18-question-evocation}{PMCE question evocation}
  {Classical propositional declaratives; finite set antecedents and ordered
  interrogative succedents; primitive negation; no equality predicate.
  All displayed lists must form questions, with at least two pairwise
  syntactically distinct answers. Definitions~\ref{def:p18-qev-sequents},
  \ref{def:p18-qev-rule-standing} and~\ref{def:p18-qev-alpha-beta}
  fix the list notation and formula types. Ax and \(R_1\)--\(R_4\)
  are primitive; later entries have the explicitly indicated standing.}

\hbrule{rule:p18-question-evocation-ax}{Ax}
  {\text{no premises}}
  {\vdash ?\{D_1,\ldots,D_n\}}
  {Each \(D_i\) is a clause without complementary literals;
   \(D_1\lor\cdots\lor D_n\) contains complementary literals.}

\hbrule{rule:p18-question-evocation-r1}{\ensuremath{R_1}}
  {X\vdash ?\{\vec A,B\}\qquad X\vdash ?\{\vec A,C\}}
  {X\vdash ?\{\vec A,B\land C\}}
  {The common list \(\vec A\) is nonempty and \(B\land C\)
   is syntactically different from every member of it.}

\hbrule{rule:p18-question-evocation-r2}{\ensuremath{R_2}}
  {X\vdash ?\{\vec A,B\}}
  {X\vdash ?\{\vec A,C\}}
  {\(B\leftrightarrow C\) is a classical tautology;
   \(\vec A\) is nonempty and \(C\) differs syntactically
   from every member of it.}

\hbrule{rule:p18-question-evocation-r3}{\ensuremath{R_3}}
  {X\vdash ?\{B\to A_1,\ldots,B\to A_n\}}
  {X,B\vdash ?\{A_1,\ldots,A_n\}}
  {\(n>1\); all displayed expressions are erotetic sequents.}

\hbrule{rule:p18-question-evocation-r4}{\ensuremath{R_4}}
  {X\vdash Q}{X\vdash Q'}
  {\(dQ=dQ'\); the question expressions themselves need not be equal.}

\hbrule{rule:p18-question-evocation-r1-position}{\ensuremath{R_1^*\ (\mathrm{derived})}}
  {X\vdash ?\{\vec L,B,\vec R\}\qquad
   X\vdash ?\{\vec L,C,\vec R\}}
  {X\vdash ?\{\vec L,B\land C,\vec R\}}
  {\(\vec L,\vec R\) together form a nonempty common list;
   \(B\land C\) differs syntactically from all its members.}

\hbrule{rule:p18-question-evocation-r2-position}{\ensuremath{R_2^*\ (\mathrm{derived})}}
  {X\vdash ?\{\vec L,B,\vec R\}}
  {X\vdash ?\{\vec L,C,\vec R\}}
  {\(B\leftrightarrow C\) is a tautology; the common list is
   nonempty and \(C\) differs syntactically from all its members.}

\hbrule{rule:p18-question-evocation-transpose}{\ensuremath{R_4^*\ (\mathrm{derived})}}
  {X\vdash ?\{A_1,\ldots,A_n\}}
  {X\vdash ?\{A_{\sigma(1)},\ldots,A_{\sigma(n)}\}}
  {\(\sigma\) exchanges two positions and fixes the others.}

\hbrule{rule:p18-question-evocation-cut-right}{\ensuremath{R_{\mathrm{cut}r}\ (\mathrm{admissible})}}
  {X\vdash ?[A\mid n,B]\qquad X\vdash ?[A\mid n,\neg B]}
  {X\vdash ?[A\mid n]}
  {\(n>1\); every premise list satisfies syntactic distinctness.}

\hbrule{rule:p18-question-evocation-cut-left}{\ensuremath{R_{\mathrm{cut}l}\ (\mathrm{admissible})}}
  {X,B\vdash ?[A\mid n]\qquad X,\neg B\vdash ?[A\mid n]}
  {X\vdash ?[A\mid n]}
  {\(n>1\).}

\hbrule{rule:p18-question-evocation-left-equivalence}{\ensuremath{R_{l\mathrm{eqv}}\ (\mathrm{admissible})}}
  {X,B\vdash ?[A\mid n]}{X,C\vdash ?[A\mid n]}
  {\(B\leftrightarrow C\) is a classical tautology.}

\hbrule{rule:p18-question-evocation-r3-reverse}{\ensuremath{R_{3r}\ (\mathrm{admissible})}}
  {X,B\vdash ?[A\mid n]}{X\vdash ?[B\to A\mid n]}
  {\(n>1\).}

\hbrule{rule:p18-question-evocation-implication-right}{\ensuremath{R_{\neg\to r}\ (\mathrm{admissible})}}
  {X\vdash ?[A\mid n\to B]}{X,\neg B\vdash ?[\neg A\mid n]}
  {\(n>1\); negation and the fixed implication suffix preserve
   syntactic distinctness by the grammar.}

\hbrule{rule:p18-question-evocation-beta-left}{\ensuremath{R_{\beta l}\ (\mathrm{admissible})}}
  {X,\beta_1\vdash ?[A\mid n]\qquad
   X,\beta_2\vdash ?[A\mid n]}
  {X,\beta\vdash ?[A\mid n]}
  {\(\beta,\beta_1,\beta_2\) are a row of
   Definition~\ref{def:p18-qev-alpha-beta}.}

\hbrule{rule:p18-question-evocation-alpha-left-contract}{\ensuremath{R_{\alpha l}^{\mathrm{c}}\ (\mathrm{admissible})}}
  {X,\alpha_1,\alpha_2\vdash ?[A\mid n]}
  {X,\alpha\vdash ?[A\mid n]}
  {\(\alpha,\alpha_1,\alpha_2\) are a row of
   Definition~\ref{def:p18-qev-alpha-beta}.}

\hbrule{rule:p18-question-evocation-alpha-left-expand}{\ensuremath{R_{\alpha l}^{\mathrm{e}}\ (\mathrm{admissible})}}
  {X,\alpha\vdash ?[A\mid n]}
  {X,\alpha_1,\alpha_2\vdash ?[A\mid n]}
  {\(\alpha,\alpha_1,\alpha_2\) are a row of
   Definition~\ref{def:p18-qev-alpha-beta}.}

\hbrule{rule:p18-question-evocation-beta-right}{\ensuremath{R_{\beta r}\ (\mathrm{admissible})}}
  {X\vdash ?[A\mid n,\beta]}
  {X\vdash ?[A\mid n,\beta_1,\beta_2]}
  {Here \(n\geq1\). Require \(\beta_1\ne\beta_2\),
   \(\beta_j\ne A_i\), and the indicated \(\beta\)-row;
   both complete answer lists must be questions.}

\hbrule{rule:p18-question-evocation-alpha-right}{\ensuremath{R_{\alpha r}\ (\mathrm{admissible})}}
  {X\vdash ?[A\mid n,\alpha_1]\qquad
   X\vdash ?[A\mid n,\alpha_2]}
  {X\vdash ?[A\mid n,\alpha]}
  {Here \(n\geq1\). Require \(\alpha\ne A_i\), the
   indicated \(\alpha\)-row, and well-formed premise questions.}

\hbrule{rule:p18-question-evocation-export}{\ensuremath{R_{\mathrm{export}}\ (\mathrm{admissible})}}
  {X,B\vdash ?[A\mid n]}{X\vdash ?[A\mid n,\neg B]}
  {\(n>1\) and \(\neg B\) differs syntactically from every \(A_i\).}
